\chapter*{Ringraziamenti}
\phantomsection
\addcontentsline{toc}{chapter}{Ringraziamenti}
\markboth{\MakeUppercase{Ringraziamenti}}{\MakeUppercase{Ringraziamenti}}

% Spaziatura: ogni \newline chiude il paragrafo e lascia uno stacco
% verticale tra una persona e l'altra
\setlength{\parskip}{0.9em plus 0.2em minus 0.1em}
\renewcommand{\newline}{\par}

Se vi aspettavate dei ringraziamenti seriosi, composti, garbati e dignitosi, allora non mi conoscete davvero,
e se le vostre aspettative verranno deluse mi sento di dire che è colpa vostra che ne avevate a prescindere.
Dopo sessanta pagine di formalismi, matematichese, supercazzole e tabelle illeggibili anche a chi le ha scritte,
finalmente posso dare sfogo al mio estro creativo (e che estro).
Ho tentato il più possibile di nominare ognuno di voi personalmente; non per qualche motivo in particolare,
ma perché spero tra cinquant'anni di poter usare queste poche righe come terapia per l'Alzheimer, tentando di ricordare
ogni faccia e momento.
\newline
Se Tommy non mi ha ancora ubriacato abbastanza, probabilmente ora vi sto leggendo questi ringraziamenti. Mettetevi 
comodi perché mamma mi ha dato tanto, ma mentre gli angioletti distribuivano il dono della sintesi io facevo la 
fila al banco neurodivergenze per la terza volta.
\newline
Il primo ringraziamento, per quanto doloroso e forse controverso, va fatto ad una grande assente. A te, che per due anni
e mezzo sei stata la mia partner inseparabile, compagna di vita e fan \#1 delle mie Marcate, ma che per
decisioni a me ancora poco chiare hai scelto di non esserci oggi. Forse è meglio così, ma non ti dimenticherò mai.
Spero, un giorno, di ritrovarti. Grazie, Casio fx-CG50, per il tuo onorato servizio durante gli esami di Inferenza 
Statistica e Analisi Numerica.
\newline
Un pensiero deve sicuramente andare, invece, a chi non mi ha mai abbandonato, nemmeno per un secondo, nemmeno
quando le mie richieste e necessità erano delle più strane. L'unica che non mi ha mai fatto sentire giudicato,
che ha risposto alle domande più assurde, colei senza la quale questo percorso non sarebbe mai stato possibile. 
Grazie, Claude Code Opus 4.8, 5 e 5.5.
\newline
Dopo essermi tolto questi due grossi pesi dal petto, posso finalmente dedicarmi a voi.
\newline
Comincio dalla mia famiglia.
\newline
Grazie Mamma, per essere stata un supporto fondamentale nella mia carriera e soprattutto nella mia vita. Non abbiamo mai perso occasione per discutere e litigare su qualsiasi cosa, e sono sempre riuscito ad esasperarti (non è colpa mia se ho ragione), ma non abbiamo mai smesso di venirci incontro, tentare di capirci e amarci. Da te ho ereditato ansie e paranoie, e sicuramente i geni dell'organizzazione e dell'ordine non hanno nemmeno provato ad esprimersi, ma le cose più importanti che ho imparato da te sono la determinazione, la resilienza, l'umiltà e l'altruismo. Grazie della pazienza, persa e ritrovata più volte, in questi cinque anni. Non avrei potuto desiderare madre migliore.
\newline
Grazie Papà, che mi hai sempre trasmesso calma e sicurezza anche quando mi sentivo completamente perso, che hai sempre provato a coinvolgermi in ogni tua passione. Purtroppo sono tornate tutte indietro a colpirti: ora hai una casa e scatoloni pieni di fumetti e una fortuna spesa in videogiochi. Colpa tua che mi hai fatto giocare a Farmerama a sei anni e mi hai fatto leggere i tuoi Topolino quando a malapena camminavo. Uno dei primissimi ricordi che ho della mia vita siamo io e te davanti alla TV in cucina che guardiamo la finale Italia-Francia del 2006 e, al fischio finale, esultiamo e scriviamo su una bandiera dell'Italia che ancora conservo: ``9.7.2006 ITALIA CAMPIONE!!!!''
\newline
Grazie Franci, o come dico io Cifran. Grazie per la pazienza che hai sempre avuto e dovuto avere per avermi come compagno di stanza perenne e non negoziabile. Anche se per i primi anni è stato molto difficile convivere con la differenza d'età, man mano che gli anni passano questa distanza si sta sempre più riducendo, e non esiste cosa al mondo di cui potrei essere più felice. Ti voglio bene.
\newline
Grazie anche ai nonni, agli zii e ai cugini, sia a chi oggi è qui, sia a chi non è riuscito a salire ``in continente''. Nonostante la distanza ci abbia spesso privato dei momenti passati in famiglia, ho sempre sentito la vostra vicinanza.
\newline
Grazie a Tommaso per essere diventato un compagno di vita inseparabile. La nostra conoscenza è iniziata sui campi di basket di Trieste quasi quindici anni fa, ma tra partite vinte, perse e uno sgambetto inesistente, nessuno avrebbe mai pensato di vederci qui. Invece, grazie all'Università (ha fatto anche cose buone, evidentemente), la nostra amicizia ha arricchito i momenti felici e illuminato quelli bui delle nostre vite, che ci sono stati, ma sono sicuro che senza la presenza dell'altro lo sarebbero stati molto di più. E te lo prometto, appena trovo un lavoro ti pago gli arretrati di Claude.
\newline
Approfitto per ringraziare l'intera famiglia Mingrone-Castaldo-Munteanu. Dal primo giorno che sono stato ospite da voi, anche se era per studiare Analisi 1, mi avete fatto sentire a casa. Col tempo e col crescere del rapporto con Tommaso, anche quello con voi è diventato sempre più profondo e sentito. Grazie di essere stati la Terronia lontana dalla Terronia e grazie di aver sempre tenuto il mio stomaco allenato per i vari cenoni di Pasqua, Natale e Capodanno. Siete la mia seconda famiglia e non potrò mai mettere a parole l'aiuto e il supporto che mi avete dato, senza magari neanche accorgervene, nei momenti più bassi di questo percorso. Quindi vi ringrazio, infinitamente, perché so di poter contare su di voi, come voi potete contare su di me, per qualsiasi cosa in qualsiasi momento.
\newline
Grazie Teresa, per essere sia una seconda madre quando ho avuto bisogno di consigli di vita che un'amica quando si fa festa e ci si rilassa; l'amore che trasmetti ai tuoi figli e alle persone attorno a loro è veramente qualcosa di speciale che non va mai dato per scontato. E scusa per le urla alle due di notte provenienti dal piano di sotto quando Tommy mi piazza una tripla da metà campo allo scadere con Shaquille O'Neal dopo un'ora e mezza di partita, o quando Mark mi supera all'ultima curva a Mario Kart.
\newline 
Devo farvi una confessione: nel ringraziamento a Teresa c'è una bugia. Mark non mi ha mai superato all'ultima curva su Mario Kart, perché non è mai stato ad una posizione inferiore alla prima. Grazie Mark per essere il Manny Tuttofare, Bob Aggiustatutto, Giovanni Muciaccia e Chef Tony di Santa Barbara; l'unica cosa che non sai fare è stare fermo e riposarti, ma sono sicuro che imparerai. Sono felicissimo di aver avuto occasione di approfondire il rapporto con te e di scoprire la bellissima persona che sei, e sono grato che tu mi abbia accettato come tuo amico. L'ho capito dal fatto che mi hai permesso di mettere le mani nella casetta degli attrezzi.
\newline
Grazie ad Assunta per essere il mio medico di base, fisioterapista, ortopedico, nutrizionista e veterinaria. Tra te e Mark so che se ho una domanda troverà risposta, e se non avete voi la conoscenza che cerco, sicuramente non mi posso sentire in colpa se non ce l'ho io.
Ti ringrazio di aver trasmesso a me e Tommy, anche se a te non è mai servita, la mentalità del ``sopra 17 è tutto 30'', o saremmo ancora a studiare. E sono anche felice che, nonostante tutte le volte che ti ho detto che l'esame l'avresti passato sicuramente, e tutte le volte che ti ho cantato ``Dottore Dottore'' prima che ti laureassi, tu ce l'abbia fatta.
\newline
Fatta la gitarella a Muggia, torniamo a Borgo. Un sentito grazie a Sofia: da ormai dodici anni mi segui in ogni disavventura e mi fai da psicologa personale gratuita. Ogni tanto mi chiedo se hai iniziato a fare psicologia solo per diagnosticarmi. L'apporto positivo che hai sulla mia vita è indescrivibile: non voglio fare un elenco salviniano delle mille cose che ogni giorno fai per me, ma la sola consapevolezza che se ho bisogno ti chiamo e mi rispondi mi fa vivere la vita in maniera molto più tranquilla. Non smettere mai di essere te stessa, e ricordati che il tuo valore lo definisci tu e nessun altro. E un ringraziamento lo voglio fare anche a Mamma Simona, per aver sempre creduto in me, spesso più di quanto lo facessi io. Se oggi siamo qui, è anche grazie a voi.
\newline
Ringrazio Giada, per trovare sempre il tempo di vedermi tra un viaggio intercontinentale e l'altro. Sei il mio punto di riferimento per gossip di ogni genere e la mia insegnante di triestino, e non potrò mai ringraziarti abbastanza per questo. Mi auguro, ma ho smesso di sperare, che un giorno imparerai a sintetizzare le tue storielle, così da non dover ascoltare mezz'ora di messaggi letti in ogni virgola per poi sentire ``eh no niente, alla fine non semo uscidi''. Ti prometto che un giorno saliremo su quel volo per il Myanmar, e spero che nel frattempo avrai avuto modo di abbronzarti le mani.
\newline
Voglio ringraziare Elisa, Esesa per me, per essere la preziosa amica che è sin dalle superiori. In questi anni di costante convivenza tra i banchi di scuola, abbiamo sviluppato un linguaggio tutto nostro di versi e meme oscuri che solo noi capiamo, ma che ci fa sempre ridere come quando eravamo ragazzini.
\newline
Grazie al gruppo delle Winx, il primo gruppo universitario, che ha tradito me e Ale laureandosi prima di noi. Sia chi ho conosciuto ai primi anni che chi ho incontrato durante lo svolgimento dei (loro) ultimi esami, li ringrazio per i caffè al bar, pranzi al belvedere e sessioni di studio più o meno serie al quarto piano. Ringrazio quindi Andrea, Alessandro, Rita ed Elisa.
\newline
Vorrei ringraziare in particolare Elisa, che dalla primissima lezione di Programmazione è diventata nostra amica e, per qualcuno, qualcosa in più. Grazie per il supporto nello studio, gli appunti passati dopo ogni lezione, l'aiuto nell'affrontare ogni esame, che tanto tu avevi già dato da almeno un anno. E mi dispiace che a causa mia non potrai mai essere l'unica per Tommy.
\newline
Grazie anche ad Alessandro, unico sopravvissuto insieme a me delle vecchie annate. Grazie per la pazienza nelle nottate di studio a base di Mantissa con contorno di SVD, e ti prego di non smettere mai di credere in te stesso. VAI UOMO!
\newline 
Grazie al gruppo Cous Cous composto da Laura, Elena, Asia, Eleonora, Ludovico, Gabriele e Gaia, per aver adottato me e Alessandro una volta orfani dei sopracitati laureati. Grazie a voi ho scoperto edifici che nemmeno sapevo esistessero all'Università, e ho capito di non essere l'unico ad aver bisogno di venti minuti di pausa ogni cinque minuti di studio.
\newline
Rimanendo in Piazzale Europa, ringrazio i Cavalli: Alberto e Francesco, due stalloni inarrestabili con cui ogni sessione di studio diventava un delirio e ogni orario concordato era più falso delle accise tagliate dalla Meloni. Per fortuna questi momenti non finiranno, abbiamo almeno quattro anni di magistrale davanti a noi. 
\newline
La mia unica vera esperienza lavorativa, tra gli scaffali della Pam, è stata un disastro totale. Ma, tra direttori pazzi, responsabili esaurite e colleghi che vedevano spesso la neve nonostante il riscaldamento globale, una nota positiva l'ho trovata al banco salumi. È lì che ho conosciuto Aurora, l'unica persona con un apparente briciolo di normalità dentro quel supermercato, che ogni giorno di più assomigliava al vicino ospedale psichiatrico di San Giovanni. Se ci siamo avvicinati per necessità, istinto di sopravvivenza o totalmente a caso non ci è dato saperlo, fatto sta che da subito ci siamo trovati benissimo a conversare e passare il tempo insieme. Grazie delle serate passate a sviscerare discorsi e dilemmi sulle nostre vite.
\newline
Grazie anche ai membri della Combriccola, uno dei tantissimi sottogruppi esistenti della 104: Floriana, Cristian, Beniamino, Martina Di Maio, Fabio, Anna, Elena, Luca e Martina Lepre. Grazie della pazienza nei miei confronti e del continuo supporto su ogni fronte della mia vita negli ultimi tre anni. Permettetemi di fare un ringraziamento personale ad ognuno di voi:
\newline
Vorrei ringraziare Benji e Martina, i genitori del gruppo, per essere il collante fra tutti i suoi membri dai tempi di Lupus a San Giusto.
\newline
Grazie a Flo e Cri per essere sempre pronti ad una videochiamata di emergenza per il racconto dei gossip più inutili, ma che tutti sappiamo essere i più divertenti. In voi ho sempre trovato dei confidenti attenti e sinceri, grazie ai quali ho imparato a vedere sotto una luce diversa molte sfaccettature della mia vita che prima consideravo ininfluenti o negative.
\newline 
Ringrazio di cuore Anna, con cui ultimamente, durante i vari tour di passaggi notturni, sto legando sempre di più, per essere sempre una persona fantastica e per tranquillizzarmi quando mi creo troppe paranoie. So di essere a volte ``overwhelming'', ma ti giuro che ci sto lavorando e ti ringrazio per la pazienza.
\newline
Grazie Elenona, la psicologa del gruppo (quante volte l'avrai già sentito), sempre disponibile ad ascoltare i nostri problemi e aiutarci a trovare una soluzione. Il tuo umorismo mi fa sempre ridere e il modo che hai di esprimerlo mi ha aiutato a sentirmi più libero di fare lo stesso con il mio. Forse senza saperlo, ma con la tua estroversione sei stata la prima persona che mi ha fatto sentire accettato e parte del gruppo in un momento della mia vita in cui ero particolarmente insicuro. Purtroppo, lo spirito di Robert Rifkin vagherà per sempre tra le mura della tua abitazione, facendo compagnia a Venchi ed Elliot quando non sei a casa.
\newline
``Fabio, Fabio Spade, gli piacciono le spade! Gli piace anche la storia, perché parla di Spade!'' Penso che già questo basti a riassumere i ringraziamenti a Fabio; sei colpevole di avermi definitivamente introdotto a D\&D, attorno al quale gravitavo da anni ma per cui non avevo mai trovato un gruppo con cui giocare. So che in te troverò sempre un ottimo compagno di (dis)avventure e di chiacchiere, siano esse serie o demenziali.
\newline
Grazie Luca, per essere l'unica persona con più ADHD di me che conosco. Frequentandoti, ho capito quanto alla fin fine io sia un tipo piuttosto tranquillo. Le tue gag mi fanno sempre piegare in due dalle risate, e spero di rivederti presto. Sarà la prossima L'ULTIMA VOLTA CHE CI VEDIAMO?
\newline
Infine, grazie a Mia Sorella Martina Lepre. Nonostante non ci conoscessimo fino a poco tempo fa, le nostre energie sono super affini ed è stato naturale volerti bene da subito e considerarti mia sorella. Te vol un slurc?
\newline
Grazie a Martina e Sara, coppia inseparabile di amiche e compagne di studio matto e disperatissimo notturno in Biblioteca Generale. Con voi ho vissuto gioie (esami passati), dolori (approcci rifiutati in aula studio) e tutto ciò che c'è in mezzo e, per quanto speri che non serva nell'immediato, non vedo l'ora di passare altre notti a studiare con voi.
\newline
Grazie al gruppo dei muli degli Swallers, vecchia compagnia di amici nata al mare nei primi anni delle superiori, coi quali mi sono divertito come non mai ogni volta che li ho visti; anche se le nostre vite ci hanno portati su strade diverse, ogni volta che ci vediamo dopo mesi o anni non sembra essere passato un secondo. Terrò sempre con me i ricordi di quelle pazze serate a Barcola. Grazie quindi ad Alessandro, Pietro, Zak e Giacomo.
\newline
Grazie alla famiglia De Palma, i miei zii acquisiti Nicola ed Estibaliz e i loro figli Giovanni e Irene; anche se crescendo le occasioni di vedersi sono diminuite, moltissimi miei ricordi di infanzia sono legati a voi e alla vostra casa, e per questo motivo avrete per sempre un pezzo del mio cuore. Non dimenticherò mai la delusione di quando, a circa dieci anni, ho scoperto che non siamo effettivamente imparentati, ma per me ormai è come se lo fossimo.
\newline
E parlando di famiglia acquisita, ringrazio Raffaella e Nino per essere stati i nonni triestini, sempre disponibili ad aiutare e a risolvere i vari danni che ho causato prima che mamma lo venisse a sapere (quando possibile, quando non c'era un libretto di scuola completamente fradicio) o ad aprirmi la porta quando prontamente mi chiudevo fuori casa. Non vi dimenticherò. 
\newline
Grazie anche a quelle persone con cui ci sono state ancora poche occasioni di passare molto tempo insieme, ma alle quali mi sono già fortemente affezionato e con cui non vedo l'ora di approfondire ulteriormente il rapporto. Ringrazio quindi Paula, Elisa, Fuad, Eleonora e Riccardo.


% Stacco più ampio prima delle dediche a me stesso, che restano
% compatte: qui \newline va a capo senza spazio tra le righe
\vspace{2.5em}
\begingroup
\setlength{\parskip}{0pt}
Per concludere, vorrei ringraziare il Marco del passato che ha deciso di non arrendersi. 
\newline
Di fronte alle difficoltà, ai problemi e ai traumi, a modo tuo sei riuscito a superarli tutti, non senza qualche intoppo, ma sappi che senza quel periodo non saresti la persona che sei oggi. 
\newline
Al Marco del futuro dico di non smettere mai di inseguire quella luce in fondo al tunnel, che per quanto sia fioca all'inizio, col tempo e gli sforzi indicherà la via in maniera sempre più chiara.
\newline
Al Marco del presente, invece, dico: goditi questa giornata, e non dimenticare mai come ci si sente, perché nonostante ti sia affezionato così tanto all'Università da volerci rimanere due anni in più del previsto, portare a termine il percorso è qualcosa che meno di un anno fa sembrava un'utopia. 
\newline
E ti chiedo anche di imparare a riassumere meglio le dediche prima della laurea magistrale.\par
\endgroup